%%%%%%%%%%%%%%%%%%%%%%%%%%%%%%%%%%%%%%%%%%%%%%%%%%%%%%%%%%%%%%%%%%%%%%%%%%%%%%%%
%%%%%%%%%%%%%%%%%%%%%%%%%%%%%%%%%%%%%%%%%%%%%%%%%%%%%%%%%%%%%%%%%%%%%%%%%%%%%%%%
% CONCLUSIONES %
%%%%%%%%%%%%%%%%%%%%%%%%%%%%%%%%%%%%%%%%%%%%%%%%%%%%%%%%%%%%%%%%%%%%%%%%%%%%%%%%

\cleardoublepage
\chapter{Conclusiones}
\label{chap:conclusiones}

En este capítulo se recogen las conclusiones obtenidas tras el desarrollo de WebBuilder. En primer lugar se evalúan los objetivos planteados al inicio del proyecto, analizando qué se ha logrado y qué limitaciones se han encontrado. A continuación se reflexiona sobre la aplicación de los conocimientos adquiridos durante el grado y las lecciones aprendidas a lo largo del desarrollo. Por último se proponen posibles líneas de trabajo futuro que podrían ampliar o mejorar el sistema en próximas iteraciones.

\section{Consecución de objetivos}
\label{sec:consecucion-objetivos}

%Y si has llegado hasta aquí, siempre es bueno pasarle el corrector ortográfico, que las erratas quedan fatal en la memoria final. Para eso, en Linux tenemos aspell, que se ejecuta de la siguiente manera desde la línea de \emph{shell}:
%
%\begin{verbatim}
%  aspell --lang=es_ES -c memoria.tex
%\end{verbatim}

Cabe destacar que una de las principales limitaciones encontradas durante el desarrollo ha sido la calidad visual de los sitios generados. La calidad que producen modelos de pago de alto rendimiento como GPT-4 o Claude no es alcanzable con los modelos gratuitos disponibles en los proveedores utilizados, y por más refinamiento aplicado en los prompts y en el sistema de presets visuales, llegar a ese nivel de calidad resulta prácticamente imposible sin acceso a modelos más potentes. Los objetivos marcados y conseguidos son los siguientes: 

\begin{itemize}
  \item \textbf{Desarrollar un módulo de ingestión de datos.} El módulo implementado en \\ \texttt{utils/ingest/} es capaz de validar la URL introducida por el usuario, descargar el contenido de la API de forma segura y detectar automáticamente el formato de los datos. Los cuatro formatos contemplados en el objetivo, JSON, XML, CSV y GeoJSON, están soportados y producen una estructura Python homogénea que el resto del sistema consume.
  
  \item \textbf{Integrar modelos de lenguaje para el análisis y la generación de código.} El módulo \texttt{utils/llm/} implementa un cliente compatible con el estándar OpenAI que permite conectar con cualquier proveedor del mercado. Los modelos de lenguaje intervienen en seis de los siete pasos del generador, produciendo desde el modelo de datos Django hasta los templates HTML y el comando de carga de datos. Se han diseñado prompts específicos para cada paso tras muchas pruebas y refinamientos.
  
  \item \textbf{Implementar un generador de proyectos Django.} El generador, implementado en \\ \texttt{utils/generator/project\_generator.py}, produce un proyecto Django completo y desplegable a partir del esquema validado por el usuario. Cada paso tiene su propio prompt, su propio fallback de emergencia y su registro en el modelo \texttt{GenerationLog}, garantizando que el resultado siempre sea desplegable incluso cuando el LLM produce una respuesta de baja calidad.
  
  \item \textbf{Automatizar el despliegue mediante Docker y n8n.} El workflow \textit{webbuilder-deploy} de n8n recibe la señal de Django mediante un webhook, construye la imagen Docker del proyecto generado, selecciona un puerto libre, arranca el contenedor y devuelve la URL de previsualización. Todo el proceso es transparente para el usuario, que recibe además una notificación por correo cuando su sitio está listo.
  
  \item \textbf{Diseñar un sistema de trazabilidad y métricas.} El modelo \texttt{GenerationLog} registra cada llamada al LLM durante la generación, incluyendo el prompt enviado, la respuesta recibida, los errores detectados y si fue necesario un reintento. El verificador de consistencia \texttt{consistency\_checker.py} detecta referencias incorrectas entre ficheros antes del despliegue. El panel de métricas, accesible para administradores, toma esta información y muestra los datos a modo de dashboard.
  
  \item \textbf{Ofrecer soporte multi-LLM y configuración personalizada.} El catálogo de modelos definido en \texttt{utils/llm/llm\_catalog.py} ofrece al usuario tres modelos preconfigurados con sus ventajas e inconvenientes. Además, el modelo \texttt{UserProfile} permite a cualquier usuario conectar su propio proveedor OpenAI-compatible introduciendo su URL base, el identificador del modelo y su clave de API, que se almacena cifrada en base de datos mediante \texttt{EncryptedCharField}.
\end{itemize}


\section{Aplicación de lo aprendido}
\label{sec:aplicacion}

Durante el desarrollo de WebBuilder se han puesto en práctica numerosos conocimientos adquiridos a lo largo del Grado en Ingeniería Telemática. A continuación se detallan las asignaturas directamente relacionadas con el proyecto:

\begin{enumerate}
  \item \textbf{Programación de Sistemas de Telecomunicación:} Esta asignatura proporcionó la base de programación en Python que ha sido el lenguaje principal del proyecto. Los conocimientos adquiridos sobre estructuras de datos, funciones, módulos y manejo de ficheros han sido esenciales para desarrollar toda la lógica del backend, desde el módulo de ingestión hasta el generador de proyectos.
  
  \item \textbf{Laboratorio de Bases de Datos:} Los conceptos de modelado relacional, consultas y gestión de datos aprendidos en esta asignatura han sido necesarios para entender como diseñar los modelos de Django.
  
  \item \textbf{Aplicaciones Telemáticas:} Esta asignatura introdujo el desarrollo web con HTML, CSS y JavaScript, tecnologías que conforman toda la capa de presentación de WebBuilder. Además, los conceptos de contenedores y entornos de ejecución aislados trabajados en esta asignatura sentaron las bases para comprender y aplicar Docker en el despliegue de los proyectos generados.
  
  \item \textbf{Servicios Telemáticos:} El framework Django, utilizado como núcleo del sistema, fue introducido en esta asignatura. El manejo de peticiones, respuestas, rutas y vistas se aplican directamente en WebBuilder.
  
  \item \textbf{Seguridad en Redes de Ordenadores:} Los conocimientos sobre seguridad y cifrado adquiridos en esta asignatura han sido aplicados directamente en la protección de las claves de API de los proveedores LLM, almacenadas cifradas en base de datos mediante \texttt{EncryptedCharField}, evitando que se guarden en texto plano.
  
  \item \textbf{Sistemas Operativos:} Aunque esta asignatura se centró en programación en C y conceptos de bajo nivel como procesos, kernel y gestión de recursos, esos fundamentos han facilitado la comprensión del funcionamiento interno de Docker: la gestión de contenedores, la asignación de puertos y el aislamiento de procesos en el sistema anfitrión.
\end{enumerate}

Este proyecto ha ayudado integrar y aplicar de forma práctica los conocimientos adquiridos a lo largo del grado, demostrando que las competencias desarrolladas en cada asignatura no son independientes sino que se complementan y convergen en un proyecto real de ingeniería.

\section{Lecciones aprendidas}
\label{sec:lecciones_aprendidas}

El desarrollo de WebBuilder ha supuesto el aprendizaje de varias tecnologías que no formaban parte de la carrera. Entre ellas destacan:

\begin{enumerate}
  \item \textbf{Automatización de flujos con n8n:} El uso de n8n para orquestar flujos de trabajo asíncronos ha sido un aprendizaje completamente nuevo. Aprender a diseñar workflows con nodos, configurar webhooks, gestionar errores en cada paso y conectar servicios externos.
   
  \item \textbf{Integración de modelos de lenguaje mediante APIs:} Conectar el sistema a proveedores de LLM como Groq, OpenRouter o NVIDIA NIM, gestionar sus límites de uso, interpretar sus respuestas y manejar sus fallos, asi como aprender a manejar el estándar OpenAI.
  
  \item \textbf{Ingeniería de prompts:} El diseño, prueba y refinamiento de prompts efectivos necesarios durante el proyecto. Aprender a formular instrucciones precisas, cómo estructurar el contexto que se le pasa al modelo y cómo proteger el sistema frente a respuestas inesperadas ha requerido un proceso largo de experimentación.
  
  \item \textbf{Tailwind CSS:} El uso de Tailwind como sistema de estilos para los proyectos generados, ha sido aprendido íntegramente durante el proyecto. Su compatibilidad natural con la generación por LLM lo han convertido en una herramienta indispensable de la generación.
  
  \item \textbf{Gestión autónoma de un proyecto de software:} Llevar un proyecto de esta envergadura de forma completamente autónoma, realizando investigación, tomando decisiones de arquitectura y añadiendo funcionalidades, también organizando y realizando refactors cuando era necesario.
  
  \item \textbf{Redacción técnica con LaTeX:} La elaboración de esta memoria ha supuesto aprender a trabajar con \LaTeX\ desde cero, incluyendo la gestión de referencias, figuras, tablas y la estructura general de un documento técnico académico.
\end{enumerate}

\section{Trabajos futuros}
\label{sec:trabajos_futuros}

Aunque WebBuilder cumple con los objetivos planteados, existen alguna lineas de mejora y ampliación. A continuación se proponen las mas relevantes, tanto mejoras técnicas del sistema como posibles funcionalidades.

\begin{enumerate}
  \item \textbf{Mejora de la calidad visual de los sitios generados:} La principal línea de trabajo futuro es mejorar la calidad visual de los proyectos generados. Esto pasa por incorporar modelos de mayor capacidad y ampliar la librería de ejemplos few-shot con más variantes de estilo.
  
  \item \textbf{Validación automática del código generado:} Actualmente el sistema detecta inconsistencias entre ficheros pero no ejecuta el proyecto antes de desplegarlo. Una mejora sería lanzar el proyecto en un contenedor efímero tras la generación, verificar que responde HTTP 200 y solo entonces marcarlo como listo, cerrando el ciclo de calidad de forma completamente automática.
  
  \item \textbf{Soporte para más formatos de entrada:} Además de URLs de APIs públicas, sería interesante permitir al usuario usar APIs o ficheros en cualquier formato, ampliando así las fuentes de datos que el sistema puede procesar.
  
  \item \textbf{Comparador de modelos LLM:} Permitir lanzar la misma generación con dos modelos distintos en paralelo y comparar los resultados sería una funcionalidad muy potente, tanto para el usuario final como para el análisis interno del sistema.
  
  \item \textbf{Galería pública de proyectos generados:} Implementar una sección donde los usuarios puedan marcar sus proyectos como públicos y compartirlos en una galería con el tipo de sitio, la API de origen y una captura visual. Convirtiendo asi WebBuilder en una herramienta social añadiendo interactividad y aumentando su potencial.
\end{enumerate}